\documentclass[12pt,a4paper]{exam}

\usepackage[utf8]{inputenc}
\usepackage[T1]{fontenc}
\usepackage[french]{babel}
\usepackage{amsmath,amssymb}
\usepackage{pgfplots}
\pgfplotsset{compat=1.18}
\usepackage{tikz}
\usetikzlibrary{arrows.meta}
\usepackage{geometry}
\geometry{margin=2cm}

\setlength{\parindent}{0pt}

\begin{document}

\begin{center}
\Large \textbf{- Devoir Surveillé -\\Suites numériques (1\textsuperscript{ère} STMG)}
\vspace{0.5cm}
\\\large \underline{\textbf{RÉPONDRE ENTIÈREMENT SUR UNE COPIE SÉPARÉE}}
\vspace{0.5cm}
\\ \textit{\underline{1 point bonus} sera accordé pour les copies soignées avec réponses 
entièrement justifiées.}
\end{center}
\vspace{0.4cm}
\noindent\hrulefill
\vspace{0.5cm}

% ---------------------------------------------------------
\textbf{Exercice 1 – Nombre de « like » (4 points)}
\vspace{5pt}

On considère le nombre de « likes » obtenus chaque semaine par un groupe de musique
sur sa page Instagram :

\[
\begin{array}{c|c|c|c|c|c|c|c|}
\text{Semaine }n & 0 & 1 & 2 & 3 & 4 & 5 & 6\\
\hline
u_n & 130 & 155 & 180 & 205 & 230 & 255 & 280
\end{array}
\]
\vspace{2pt}
\begin{questions}
    \question[1,5] Quelle est la nature de \((u_n)\) ? Donner sa raison.
    \question[1,5] Donner l’expression explicite de \(u_n\) en fonction de $n$.
    \question[1] Sachant qu'une année comporte 52 semaines, calculer le nombre de « likes »
     au bout d'un an.
\end{questions}

\noindent\hrulefill
\vspace{0.5cm}

% ---------------------------------------------------------
\textbf{Exercice 2 – Disneyland Paris (6 points)}
\vspace{5pt}

Un abonnement illimité à Disneyland Paris coûtait 470 € en 2023.  
Chaque année, son prix augmente de 20 €.  
On note \(p_n\) le prix \(n\) années après 2023.

\begin{questions}
    \question[1] Expliquer pourquoi la situation modélise une suite arithmétique.
    \question[1] Donner \(p_0\) et la raison \(r\).
    \question[1] Exprimer $p_{n+1}$ en fonction de $p_n$.
    \question[2] Donner la formule explicite de \(p_n\) et calculer le prix en l'an 2033.
    \question[1] À partir de quelle année dépassera-t-on 800 € ?
\end{questions}

\noindent\hrulefill
\vspace{0.5cm}

% ---------------------------------------------------------
\textbf{Exercice 3 – Modèle géométrique : remises successives (6 points)}
\vspace{5pt}

Un produit en solde ne se vend pas. Chaque jour, le vendeur baisse son prix de 10\%.
Son prix initial est de 30 €.

On note \(u_n\) le prix de l'article le \(n\)-ième jour.

\begin{questions}
    \question[1] Donner $u_0$ et calculer $u_1$, $u_2$ et $u_3$.
    \question[1,5] Exprimer $u_{n+1}$ en fonction de $u_n$.
    \question[1] Exprimer \(u_n\) en fonction de $n$.
    \question[1] Calculer $u_{30}$. On donnera un arrondi au dixième.
    \question[1,5] Le prix sera-t-il inférieur à 10€ au bout de 10 jours ? Justifier.
\end{questions}

\noindent\hrulefill
\vspace{0.5cm}

% ---------------------------------------------------------
\textbf{Exercice 4 – Étude d'une suite (4 points)}
\vspace{5pt}

On considère la suite $(u_n)$ définie pour tout entier naturel $n$ par :
\[
u_n = n(n+2) + 1.
\]

\begin{questions}
    \question[1] Calculer les 3 premiers termes de la suite.

    \question[1] Exprimer $u_{n+1}$ en fonction de $n$.

    \question[2] Montrer que $(u_n)$ est croissante. Justifier.
\end{questions}

\noindent\hrulefill


\vspace{0.5cm}
\begin{center}
\textbf{Fin du devoir — relire la copie avant de rendre.}
\end{center}

\end{document}
